\begin{frame}{4. Transparencias durante el trabajo paralelo}
  En la práctica utilizamos \texttt{parallelize()} para distribuir cuatro vehiculos en dos y cuatro particiones.

  \begin{itemize}
    \item Transparencia ofrecida: Spark distribuye los datos y genera las tareas sin que tengamos que asignar manualmente cada vehiculo.
    \item Transparencia no completa: la distribucion y ejecucion interna no se muestran en la salida normal del programa. Para observarlas utilizamos Spark UI.
  \end{itemize}

  Comprobamos las particiones, los Jobs, los Stages y las Tasks desde la interfaz web.

  \par
\end{frame}

\begin{frame}{5. Comunicacion que coordina tareas}
  La practica se ejecuto en modo local mediante \texttt{local[*]}.

  \begin{itemize}
    \item Spark mostro un executor identificado como \texttt{driver}, con 12 cores disponibles.
    \item Al utilizar dos particiones observamos dos tareas completadas.
    \item Con cuatro particiones observamos cuatro tareas completadas.
    \item En el Event Timeline las tareas comenzaron aproximadamente al mismo tiempo.
  \end{itemize}

  La practica permitio observar la coordinacion de tareas en modo local, pero no comprobo la comunicacion entre nodos de un cluster ni un protocolo de mensajes.

  \par
\end{frame}

\begin{frame}{6. Espera y avance de tareas}
  Para recuperar los resultados utilizamos la accion \texttt{collect()}, que espera la ejecucion de las operaciones pendientes y devuelve los resultados al programa principal.

  \begin{itemize}
    \item Las tareas finalizaron correctamente: 2 de 2 y 4 de 4.
    \item El Event Timeline mostro intervalos de ejecucion superpuestos.
    \item No implementamos acciones asincronas ni barreras de sincronizacion.
  \end{itemize}

  Tambien medimos el tiempo de ejecucion. Con dos particiones obtuvimos un promedio de 1.8195 segundos y con cuatro particiones 3.4206 segundos.

  \par
\end{frame}
